\chapter{总结与展望}
\section{论文工作总结}
本论文面向工程化检索与数据采集场景，设计并实现了一套以 Search Agent 为核心的端到端检索流水线，旨在将用户自然语言提示通过模型驱动的解析，工程化地对接实时在线与离线垂域检索后端，完成检索、可选的实时爬取、解析、清洗、质量评估与索引入库的全流程。论文的工作可从方法论贡献与工程实现两大层面进行概括。

在方法论层面，本论文系统化地提出并实现了 prompt-to-query 的工程化路线：结合规则与轻量模型的混合解析架构，输出可解释、带置信度的 QueryPackage，包括关键词、意图标签、触发爬取标识等，并通过在线轻量化推理满足交互延迟约束，同时为离线强模型提供回流校正样本。针对混合检索问题，本论文阐述了并实现倒排检索与向量检索并行召回、跨通道 score normalization 与多阶段融合策略，提出在候选层面采用两阶段去重（精确 id 去重与近似相似度去重）与配置化通道配额，以在保证语义覆盖的同时控制冗余。重排方面，论文采用分层重排思想：以 feature-based 的轻量 LTR 保证在线延迟需求，必要时对 top-K 启用 cross-encoder 做精排，并通过蒸馏与特征缓存减小二阶段模型对实时性能的影响。此外，论文在实时爬取与解析方面提出了“先快后补”的工程策略：在判断需要页面详情时异步触发爬取，不阻塞首屏响应，并设计了域名级限速、渲染降级、严格超时与断点续传策略，以平衡新鲜度、响应性与外部合规性。

在工程实现层面，论文构建了端到端的消息驱动流水线与模块化服务框架：以 MCP 为协调层实现同步与异步任务管理、路由与回调；Search Agent 作为决策层实现 prompt 路由、触发爬取与候选融合；Prompt Parser（在线与离线）提供语义解析服务；实时适配器与爬虫池完成低延迟召回与页面抓取；数据处理模块负责实时解析、二次清洗与质量评估；索引层分别实现倒排索引（文本检索）与向量索引（语义检索）的增量与批量建库策略。为保证系统的工程可用性，论文还讨论了监控、日志、告警、审计、灰度发布与回滚等运维措施，并提出了具体的存储与索引策略（对象存储保存原始抓取、关系数据库保存元数据、倒排与向量索引做检索支撑）以支持回放重跑与审计。

在评估与验证方面，论文制定了完整的指标体系：包含系统性能指标（P50/P95/P99 延迟、吞吐与 QPS、资源利用率）、检索质量指标（Recall@K、nDCG、MRR）、数据工程指标（抓取成功率、解析错误率、抓取到可检索延迟）及用户感知指标（CTR、满意度调查等），并以离线回放、A/B 小流量灰度与压力测试等方式进行验证。总体实验与工程验证表明：模型驱动的 prompt-to-query 能有效提高检索的语义相关性；实时爬取在受控超时与并发策略下可以显著提升结果的新鲜度；跨通道融合与分层重排在延迟约束下提升了最终排序效果。同时论文也识别出系统在索引更新成本、实时爬取稳定性、模型输出解释性等方面的局限与挑战。

最后，论文在实现细节上给出了可部署性的最佳实践：包括解析模块的版本化与灰度发布流程、爬取池的资源隔离策略、索引构建的 alias 原子切换、消息队列驱动的异步补全以及完整的监控、告警与审计链路。这些工程化细节使得提出的方法不仅具备学术价值，也具有明确的工业落地潜力。
\section{后续工作展望}
尽管本论文实现了一个功能完备且工程化的 Search Agent 系统，但在若干重要方向上仍存在可延展和优化的空间。以下从技术深化、系统工程、合规与应用扩展等角度阐述后续工作方向：

1.未来可在Prompt 解析与语义理解方面引入在线学习和微调机制，提升对新兴实体和用户表达的适应能力。此外，实时爬取可通过优化反爬机制与渲染策略，提升资源利用效率。在向量索引与增量更新上，研究高效的在线索引结构和增量更新算法，以减少重建成本。同时，在召回融合与重排方面，探索端到端学习框架和在线学习能力，提升系统的灵活性和实时性。

2.加强隐私保护与合规审查，引入差分隐私等技术确保数据安全，并完善自动化合规机制。未来还将扩展多模态检索与生成模型集成，增强系统在不同场景下的适应性。此外，通过构建成熟的MLOps体系，实现模型与数据版本治理、自动化回放与性能回归检测，以支撑大规模生产环境中的稳定运行。